% -*- coding: utf-8 -*-
\chapter{总结与展望}
\label{chap:conclusion}

\section{全文工作总结}
\label{sec:summary}

低空智联网作为支撑低空经济安全高效运转的数字中枢，对低空目标的实时连续感知与宽带数据交互提出了严苛的协同要求。本文针对低空无人机系统面临的空地强多径衰落、多普勒扩展以及通感资源竞争等关键矛盾，聚焦于一个边界清晰的核心问题：在通信需求被完全满足的前提下，如何把剩余资源最优地用于感知，并给出可横向比较的公平评价。围绕这一问题，本文完成了从系统建模、理论推导、算法构造到仿真验证的完整链路，主要成果可归纳为以下五项。

\textbf{第一，建立了面向低空立体空域的空地多载波通感一体化系统模型。} 模型统筹考虑了无人机三维几何拓扑、莱斯衰落下行信道、双程雷达回波与阵列空间导向矢量四个要素，并给出了发射协方差矩阵对通信速率与感知精度双重映射关系的显式表达。模型的一个关键设计是在接收端引入匹配空域滤波器，利用目标方向与用户方向的空间分离实现回波与通信信号的空域解耦，从而使同一时频资源上的通感共存具备物理可实现性。

\textbf{第二，推导并验证了角度估计 CRB 的闭式表达式，纠正了阵列标度的系统性误用。} 本文证明在回波幅度吸收全部阵列增益的建模约定下，均匀线阵角度估计的 CRB 为
\begin{equation}
  \mathrm{CRB} = \frac{6\sigma^{2}}{\pi^{2}\cos^{2}\theta\,\bigl(N_{\mathrm{R}}^{2}-1\bigr)\lvert\alpha\rvert^{2}N_{\mathrm{sym}}\,G(\theta)},
  \label{eq:concl_crb}
\end{equation}
其随阵元数的标度为 $N_{\mathrm{R}}^{-2}$，而非文献中广泛引用的 $N^{-3}$。二者的差异源于信噪比定义的前提不同：三次方标度假设信噪比按单阵元定义，而本文的 $\lvert\alpha\rvert^{2}$ 已包含全部阵列增益。这一纠正并非文字之争，它直接影响资源分配决策：若误用三次方标度，会系统性高估增加阵元的感知收益。第四章以五个阵元数配置作蒙特卡洛验证，实测均方根误差与闭式解的比值落在 $[0.938,\,1.050]$ 区间内，且 CRB 随阵元数加倍而降至约四分之一，与理论标度一致。

\textbf{第三，提出等通信速率约束下最小化感知 CRB 的建模方式，替代了加权和目标。} 本文论证了加权和目标的物理不可解释性：两个分量量纲不同（bit/s/Hz 与 rad$^{2}$），权重缺乏物理含义；更严重的是权重趋于零时系统退化为单一功能优化，另一功能的性能可发散至非物理奇点。据此本文改以“通信速率不低于门限”为硬约束、以最小化感知 CRB 为目标，使优化结果具有明确的工程解释：在保证业务速率的前提下，感知精度能达到的极限。

\textbf{第四，构造了可证明最优性的低复杂度求解方案。} 本文放弃需要秩一恢复的半正定松弛路线，转而直接构造满足全部约束的可行解：以迫零或 MMSE 确定通信波束方向、以匹配滤波确定感知波束、以一维闭式表达确定功率分配。本文证明了两条结构性质：其一，由于感知增益关于各功率变量线性且系数非负，“通信功率取可行下界、剩余功率全部用于感知”是最优结构；其二，在窄带假设下各子载波对感知增益贡献相同，故按“单位功率通信效率”排序选取感知子载波是全局最优的，交换论证将规模为 $\binom{M}{n_s}$ 的组合搜索降为 $O(M\log M)$ 的一次排序。此外本文给出了临界目标速率的定义，说明功率受限系统存在硬性速率上限，这解释了第四章 Pareto 曲线为何止于 $\gamma=6$\,dB。

\textbf{第五，建立了“等 CRB 速率增益”的公平对比框架并给出可复现的诚实评测规范。} 该框架以相同感知 CRB 为基准比较各方案可支撑的通信速率，避免了加权和框架下不同文献结果无法横向对照的问题。基于该框架的实测结论是：联合优化相对时隙划分在等速率下的感知优势为 1.28$\times$–1.94$\times$，且随目标速率提高单调扩大；但该优势并非处处成立，在 $k=48$ 处因感知子载波与通信子载波出现结构性重合，增益降为 0.00\%，$k=63$ 处回升至 4.95\%。本文如实报告了这两个非单调点，而非只挑选有利区间。同时，全部仿真脚本、参数配置与结果数据随论文一并提供，对不可行点、非单调点、轨迹初始条件不对称以及未实现模块逐一如实标注。

\section{方法的适用边界与局限}
\label{sec:boundary}

学术研究的价值不仅在于结论本身，也在于明确结论成立的边界。本文方法的适用边界可归纳为六条，均已在第四章相应小节中标注。

\textbf{其一，非全局最优。} 本文构造的是可行解而非 $\mathcal{P}_{1}$ 的全局最优解。第三章已论证问题的非凸性来自速率约束的对数行列式结构与 CRB 对波束的双线性依赖两处，全局最优需借助 WMMSE 类迭代或半正定松弛加秩一恢复，而后者会破坏速率约束可行性。本文的取舍是：牺牲全局最优性保证，换取解的可行性、可复现性与物理可解释性。

\textbf{其二，单目标单径。} 感知模型假设单个点目标、单条直射回波，未考虑多目标、多径与扩展目标场景。多目标时 CRB 需扩展为费舍尔信息矩阵的迹或行列式，多径时需处理相干源导致的秩亏问题，二者都会改变优化结构。

\textbf{其三，理想信道状态信息。} 本文假设基站已知全部用户的瞬时信道与目标的精确角度。实际系统中信道估计误差与角度量化误差会同时恶化速率约束的可行性与感知增益的准确性。鲁棒优化需在约束中引入误差球或采用机会约束，其代价是可行域收缩。

\textbf{其四，未建模有限字符与射频损伤。} 本文的速率表达基于高斯信令假设，未考虑有限字符星座、功率放大器非线性、相位噪声与量化误差。这些因素在高峰均比场景下会显著降低实际可达速率，也会破坏感知波形的模糊函数理想性。

\textbf{其五，轨迹对比的初始条件不对称。} 本文的轨迹实验如实显示：优化轨迹在第 1 时隙的感知精度反而劣于盘旋基线，原因是优化轨迹起点距目标 412\,m 与 654\,m，而基线起点距目标仅 269\,m。优势只在进入稳态（约第 20 时隙之后）后显现，稳态 CRB 比值约 52 倍。这一结果被保留而非隐藏，因为它揭示了轨迹优化的一类固有性质：优化收益依赖初始条件与优化时域长度，短时域下未必优于设计良好的固定轨迹。

\textbf{其六，未对比学习类方法。} 本文未引入深度强化学习等学习类基线。学习类方法在轨迹优化上已有报道，但其可解释性与可复现性较弱、训练成本高，且与本文的凸优化框架不属同一技术路线，强行对比难以保证公平。这一空白留待后续工作。

\section{未来研究方向展望}
\label{sec:future}

结合未来 6G 网络演进趋势与低空产业实际发展需求，后续研究工作可在以下七个方向进一步拓展。

\textbf{其一，宽带斜视条件下的子载波选择修正。} 本文的交换论证依赖“各子载波感知贡献相同”这一窄带前提。当带宽达到数吉赫兹时，载频在不同子载波上的相对偏移不可忽略，导向矢量随子载波变化而产生波束斜视效应，排序准则需相应修正为依赖子载波频率的加权排序。这一扩展不改变交换论证的框架，但会改变排序键。

\textbf{其二，多目标与多径场景的感知度量扩展。} 多目标时 CRB 应扩展为费舍尔信息矩阵的迹或行列式，其关于功率的线性性不再成立，闭式功率分配需要重新推导。多径场景下相干源会导致费舍尔信息矩阵秩亏，需引入解相干处理或改用子空间类算法的性能界。

\textbf{其三，非理想 CSI 下的鲁棒优化。} 可将信道误差建模为范数球或统计分布，分别采用最坏情况鲁棒优化与机会约束优化。前者保守但保证可行，后者在给定中断概率下更紧致。两种路线都会使原问题的复杂度上升一个量级。

\textbf{其四，有限字符与射频损伤的联合建模。} 可将有限字符星座的互信息下界替代高斯速率表达，并将功放非线性与相位噪声建模为等效噪声项，从而在优化中显式体现峰均比与可达速率的耦合。这一方向与波形层面的低 PAPR 设计形成互补。

\textbf{其五，半实物与实测验证。} 本文全部结论基于数值仿真。后续可结合低空飞行测试平台采集的三维空间信道冲激响应与实测雷达散射截面积，在软件无线电平台上完成算法的半实物验证，并进一步开展外场试验网评测。实测数据的引入也会检验本文理想化假设的偏离程度。

\textbf{其六，学习类基线与凸优化框架的融合。} 可将深度强化学习用于轨迹的粗粒度决策，再用本文的凸优化框架作细粒度功率与子载波分配，形成“学习定框架、优化定细节”的分层结构。这一混合路线有望同时获得学习方法的实时性与优化方法的可解释性。

\textbf{其七，多小区干扰与多视角协同感知的扩展。} 本文始终假设单个基站服务自身覆盖范围内的用户与目标，忽略了邻区干扰。在实际低空组网中，相邻基站的波束会同时照射到同一空域，通信侧需把小区间干扰纳入速率约束，感知侧则可利用多基站从不同视角观测同一目标，以多视几何换取更高的定位精度。这一扩展会使原问题的耦合维度从“子载波与用户”上升到“子载波、用户与小区”三层，分布式求解与有限信息交互将成为必要手段；同时，多视角观测下各站回波的信噪比差异较大，如何在融合时对各视角赋权以避免弱视角拖累整体精度，也是一个值得研究的问题。

\section{结语}
\label{sec:closing}

本文的全部工作可以概括为一句话：把“通感权衡”这个含混的工程直觉，转化为一个约束明确、可证明、可复现、且边界被如实标注的优化问题。这一转化的价值不在于给出了某个惊人的增益数字，而在于建立了一套可被他人检验的方法论：指标必须有物理含义，标度必须有推导依据，结论必须标注适用边界，不可行点与不利结果必须如实报告。

低空智联网的工程落地还有很长的路要走，本文所解决的只是其中资源分配环节的一个子问题。但本文所坚持的方法论立场——宁可报告一个诚实的局部结论，也不报告一个无法复现的全局承诺——或许比具体结论更具持久价值。
